\title{Click: Controllable Text Generation with Sequence Likelihood Contrastive Learning}

\begin{abstract}
It has always been an important yet challenging problem to control language models to avoid generating texts with undesirable attributes, such as toxic language and unnatural repetition. We introduce \textsc{Click} for controllable text generation, which needs no modification to the model architecture and facilitates out-of-the-box use of trained models. It employs a contrastive loss on sequence likelihood, which fundamentally decreases the generation probability of negative samples (i.e., generations with undesirable attributes). It also adopts a novel likelihood ranking-based strategy to construct contrastive samples from model generations. On the tasks of language detoxification, sentiment steering, and repetition reduction, we show that \textsc{Click} outperforms strong baselines of controllable text generation and demonstrate the superiority of \textsc{Click}'s sample construction strategy.\footnote{The project repository is available at \url{https://github.com/chujiezheng/Click}.}
\end{abstract}

\section{Introduction}

Current language models trained on massive textual corpora have shown the impressive capability of generating fluent and grammatical text \cite{gpt2, gpt3, blenderbot}. However, they often produce behaviors misaligned with human expectations. For instance, language models may generate offensive language or agree with toxic input \cite{bad, realtoxicityprompts, diasafety}. They may also generate text with unnatural repetition \cite{topp, simctg}, which is a notorious issue in autoregressive language generation. Controlling language models to avoid such undesirable attributes has always been an important yet challenging problem in NLG research.

As a popular practice, growing recent work has investigated how to decrease the generation probability of these negative samples (i.e., generations of undesirable attributes). For instance, Unlikelihood Training \cite{unlikelihood} minimizes the likelihood of each token in negative samples. GeDi \cite{gedi}, DExperts \cite{dexperts}, and Director \cite{director} adjust the next-token prediction distribution at each generation step to avoid token choices that would potentially lead to undesirable attributes.

In this work, we introduce \textsc{Click}, a method for \textbf{C}ontrollable text generation with sequence \textbf{Li}kelihood \textbf{C(K)}ontrastive learning. It employs a max-margin contrastive loss on sequence likelihood in addition to standard language modeling (\S~\ref{subsec:contrastive}), which fundamentally reduces the probability of a negative sample being decoded. Compared with previous methods of controllable text generation, \textsc{Click} has two unique advantages.
\textbf{First}, \textsc{Click} contrasts the \textit{sequence likelihoods} of positive and negative samples with a \textit{maximum likelihood margin}, which enables a higher degree of freedom for optimization than explicitly minimizing the \textit{likelihood of each token} of negative samples.
\textbf{Second}, \textsc{Click} needs \textit{no modification to the model architecture} and thus does not require laborious adjustments to the next-token prediction distribution during generation, which makes it convenient for out-of-the-box use of trained models.

We also design a likelihood ranking-based strategy of contrastive sample construction for \textsc{Click} (\S~\ref{subsec:construction}).
Given an input prompt, \textsc{Click} first samples multiple generations from the initial language model, which are labeled as \colorbox{positive}{positive}/\colorbox{negative}{negative} by a label function. It then pairs each negative sample with the positive one whose likelihood ranks highest but lower than the former. For instance, in Figure~\ref{fig:overview}, \colorbox{negative}{negative rank 2} is paired with \colorbox{positive}{positive rank 4} to constitute a pair of contrastive samples. This strategy derives from our two intuitions.
\textbf{First}, a high-likelihood positive sample (e.g., \colorbox{positive}{positive rank 1}) does not necessitate further enlargement of its likelihood gap with the negative one, which may instead result in overfitting the positive sample.
\textbf{Second}, sequence likelihood indicates how much a text is probable to be the continuation of the input, which somewhat reflects the quality of generated continuations, such as fluency and coherence. A pair of samples with a too large likelihood gap (e.g., \colorbox{negative}{negative rank 2} and \colorbox{positive}{positive rank 6}) may thus bias contrastive learning toward other aspects (e.g., fluency or coherence) than the attributes we aim to control.

We experiment with three controllable text generation tasks: language detoxification, sentiment steering, and repetition reduction (\S~\ref{sec:experiment}).
Through both automatic and human evaluation, we show that \textsc{Click} can effectively avoid undesirable attributes and outperform strong baselines. Ablation analysis further proves the superiority of \textsc{Click}'s sample construction strategy.

\section{Methodology}
\label{sec:method}

\subsection{Task Formulation}
\label{subsec:formulation}

Given an input text $x$ as a prompt, the task of controllable text generation aims to generate a fluent natural language continuation $y$ that avoids an undesirable attribute (e.g., toxicity) while maintaining contextual coherence. We denote the language model parameterized by $\theta$ as $P_{\theta}$, which produces $y$ given $x$ following the distribution $P_{\theta} (\cdot | x)$. Following the setting of controllable text generation \cite{dexperts, quark}, we also assume a label function $c(x, y)$ that assigns a binary attribute label $0/1$ to each $(x, y)$ pair\footnote{
  We assume the label function with \textit{binary} outputs rather than \textit{continuous} outputs (e.g., from 0 to 1) due to two considerations. (1) Since the label function is usually implemented as an automatic classifier, its continuous output score may be imperfect, as discussed in \S~\ref{sec:limitation}. 
  Optimization toward continuous scores may inherit more biases from the classifier, which can be alleviated to some extent by transforming continuous scores into binary labels. (2) This setting can be naturally generalized when the label function is human annotators \cite{instructgpt}, where only binary or discrete labels can be obtained. }, corresponding to a negative/positive sample, respectively.

\subsection{Sequence Likelihood Contrastive Learning}
\label{subsec:contrastive}

\textsc{Click} adopts a contrastive loss on sequence likelihood, which trains the model to assign lower generation probabilities to negative samples than positive ones. It does not need any modification to the model architecture, which makes it convenient for out-of-the-box use.
Figure~\ref{fig:overview} gives the overview of \textsc{Click}.
We first introduce how \textsc{Click} trains the language model to avoid undesirable behaviors (the 3rd step in Figure~\ref{fig:overview}) and later describe \textsc{Click}'s strategy of constructing contrastive samples in \S~\ref{subsec:construction} (the 1st and 2nd steps in Figure~\ref{fig:overview}).

\textsc{Click} requires two training sets. The first is the \textbf{language modeling set} $\mathcal{D}_\mathrm{LM} = \left\{ (x_i, y_i) \right\}_{i}$, which by default contains only positive samples, i.e., $c(x, y) = 1, \forall (x, y) \in \mathcal{D}_\mathrm{LM}$.
\textsc{Click} performs standard language modeling on $\mathcal{D}_\mathrm{LM}$ using the conventional negative log-likelihood loss:

\begin{align} \label{equ:lm}
    \mathcal{L}_\mathrm{LM} = \mathbb{E}_{(x, y) \sim \mathcal{D}_\mathrm{LM}} [ - \log P_{\theta} (y | x) ].
\end{align}

The second training set is the \textbf{contrastive learning set} $\mathcal{D}_\mathrm{CL} = \left\{ (x_i, \hat{y}_i^{+}, \hat{y}_i^{-}) \right\}_i$, which contains model-generated positive-negative sample pairs, where $c(x, \hat{y}^{+}) = 1 \wedge c(x, \hat{y}^{-}) = 0, \forall (x, \hat{y}^{+}, \hat{y}^{-}) \in \mathcal{D}_\mathrm{CL}$.
Note that the same prompt $x$ could be shared by multiple triples in $\mathcal{D}_\mathrm{CL}$.
\textsc{Click} then performs contrastive learning on $\mathcal{D}_\mathrm{CL}$ via a max-margin contrastive loss on sequence likelihood $P_{\theta} (y | x)$:

\begin{align} \label{equ:cl}
\begin{aligned}
    \mathcal{L}_\mathrm{CL} =&\ \mathbb{E}_{( x, \hat{y}^{+}, \hat{y}^{-} ) \sim \mathcal{D}_\mathrm{CL}} [ \max (0, \\
    &\gamma + \log P_{\theta} (\hat{y}^{-} | x) - \log P_{\theta} (\hat{y}^{+} | x) ) ],
\end{aligned}
\end{align}

where $\gamma$ is the margin hyperparameter. The overall optimization objective is the summation of the above two losses:

\begin{align} \label{equ:all}
  \mathcal{L} = \mathcal{L}_\mathrm{LM} + \alpha \mathcal{L}_\mathrm{CL}, 
\end{align}

where $\alpha$ is the weight hyperparameter. With the participation of $\mathcal{L}_\mathrm{CL}$, the information from the label function $c$ is injected into sequence likelihood given by the language model $P_{\theta}$.
It thus learns to avoid undesirable attributes by decreasing the generation probability of negative samples. From another perspective, we can view $\mathcal{L}_\mathrm{LM}$ as a regularization item, which maintains the language model's underlying capability of language generation.

\subsection{Contrastive Sample Construction}
\label{subsec:construction}

Before training the language model, \textsc{Click} first constructs the contrastive learning set $\mathcal{D}_\mathrm{CL}$.
We first present the overall procedure and then elaborate on the details of sample construction.

\paragraph{Overall Procedure} We start from a prompt set $\mathcal{D}_\mathrm{Pmt} = \left\{ x_i \right\}_i$, which can be easily obtained from $\mathcal{D}_\mathrm{LM}$.
For each prompt $x$ in $\mathcal{D}_\mathrm{Pmt}$, we sample multiple continuations with $P_{\theta} (\cdot | x )$, which can be implemented with popular sampling-based decoding algorithms like nucleus sampling \cite{topp}. Using the label function $c(x , \cdot )$, we split the model-generated continuations into the positive and negative sample sets $\widehat{\mathcal{Y}}^{+} = \left\{ \hat{y}_k^{+} \right\}_{k}$ and $\widehat{\mathcal{Y}}^{-} = \left\{ \hat{y}_k^{-} \right\}_{k}$.
Note that in many cases, there are fewer negative samples than positive ones, such as toxic language \cite{bad, red-teaming}. We thus pair each\footnote{
  In practice, due to the limitation of computational resources and efficiency, for each prompt $x$ we constructed at most $k$ pairs of contrastive samples for $\mathcal{D}_\mathrm{CL}$, where $k$ varies from tasks in our experiments. } negative sample $\hat{y}^{-} \in \widehat{\mathcal{Y}}^{-}$ with a positive one $\hat{y}^{+}$ and add $(x, \hat{y}^{+}, \hat{y}^{-})$ into $\mathcal{D}_\mathrm{CL}$.

\paragraph{Motivation}
A straightforward practice of sample pairing is random sampling, i.e., we randomly sample a $\hat{y}^{+} \in \widehat{\mathcal{Y}}^{+}$ for each $\hat{y}^{-}$.
However, we argue that such a practice is \textit{suboptimal}.
For a $\hat{y}^{+}$ that already has a higher likelihood than $\hat{y}^{-}$, it is unnecessary to further enlarge their likelihood gap. On the other hand, likelihood indicates how much a text is probable to be the continuation of a prompt, which somewhat reflects the quality of generated continuations like fluency and coherence. If we use $\hat{y}^{+}$ with much lower likelihoods than $\hat{y}^{-}$ to construct $\mathcal{D}_\mathrm{CL}$, contrastive learning may be biased toward other aspects (e.g., fluency or coherence) than the attributes we aim to control. Meanwhile, Equation \ref{equ:cl} would also implicitly increase the generation probability of potentially low-quality $\hat{y}^{+}$ (with low likelihoods), which conflicts with the language modeling objective (Equation \ref{equ:lm}) and may thus impair the language generation capability.

\paragraph{Likelihood Ranking-Based Strategy} Based on the above intuitions, \textsc{Click} adopts a novel likelihood ranking-based strategy for constructing contrastive samples. From the $\hat{y}^{+}$ with lower likelihoods than $\hat{y}^{-}$, \textsc{Click} selects the highest-ranked $\hat{y}^{+}$.
With the positive and negative samples at a similar likelihood level, it enables contrastive learning to focus better on the controlled attributes and also alleviates the conflict with the language modeling objective. The strategy is formulated as follows:

\begin{align} \label{equ:pair}
    \mathop{\arg \max}_{ \left\{ \hat{y}^{+} \in \widehat{\mathcal{Y}}^{+} | P_{\theta} (\hat{y}^{+} | x) < P_{\theta} (\hat{y}^{-} | x) \right\}} P_{\theta} (\hat{y}^{+} | x).
\end{align}

If all the $\hat{y}^{+}$ have lower likelihoods than $\hat{y}^{-}$, Equation \ref{equ:pair} degenerates to selecting the positive sample with the lowest likelihood, i.e., $\mathop{\arg \min}_{\widehat{\mathcal{Y}}^{+}} P_{\theta} (\hat{y}^{+} | x)$. The 2nd step in Figure~\ref{fig:overview} illustrates how our construction strategy works, where three pairs of contrastive samples are constructed: \colorbox{negative}{2}/\colorbox{positive}{4}, \colorbox{negative}{3}/\colorbox{positive}{4}, and \colorbox{negative}{5}/\colorbox{positive}{6}.

\subsection{Relationship to Prior Work}
\label{subsec:prior}

\textsc{Click} builds upon two disjoint ideas from previous work in controllable or conditional text generation.

(1) Inspired by Unlikelihood Training \cite{unlikelihood}, \textsc{Click} trains the language model to decrease the generation probability of negative samples (Equation \ref{equ:cl}).
However, Unlikelihood Training minimizes the likelihood of each token given the prefix of the negative sample, which is a \textit{token-level} objective. Different from it, \textsc{Click} adopts a max-margin contrastive loss at the \textit{sequence level}. By directly acting on sequence likelihood and setting a maximum margin $\gamma$, \textsc{Click} allows a higher degree of freedom for optimization (e.g., focusing on certain tokens that lead to undesirable attributes).

(2) Inspired by BRIO \cite{brio} and SLiC \cite{slic}, \textsc{Click} employs the contrastive loss directly on sequence likelihood. However, BRIO and SLiC align sequence likelihood with the similarity to reference text, which is not applicable for controllable text generation tasks where reference texts are usually unavailable and generation is open-ended. Unlike them, \textsc{Click} aligns sequence likelihood with the controlled attribute (the undesirable attribute corresponds to lower likelihood). Furthermore, the contrastive samples in BRIO and SLiC are randomly paired, while \textsc{Click} is based on likelihood ranking, which provides more insights about and is more tailored for open-ended text generation tasks, as verified in \S~\ref{subsec:ablation}.

\section{Experiments}
\label{sec:experiment}

We next show that \textsc{Click} can effectively avoid undesirable attributes on three controllable text generation tasks: (1) language detoxification (\S~\ref{subsec:toxicity}), (2) sentiment steering (\S~\ref{subsec:sentiment}), and (3) repetition reduction (\S~\ref{subsec:repetition}).
We also conduct ablation analysis to give further insights about \textsc{Click} (\S~\ref{subsec:ablation}).

\subsection{Language Detoxification}
\label{subsec:toxicity}

Language models are known to produce offensive language \cite{realtoxicityprompts} or express agreement with toxic input \cite{bad, diasafety}, which potentially hinders downstream tasks and real-world applications \cite{red-teaming, augesc}. The task of language detoxification aims to avoid toxic and unsafe generations.

\paragraph{Experimental Setups} We evaluated on the Bot-Adversarial Dialogue (BAD) \cite{bad} dataset. It contains human-bot conversations where human adversarially induces language models to produce unsafe generations. Each utterance is annotated with binary labels (safe or unsafe). We use the official data split, see Appendix~\ref{sec:dataset} for dataset statistics. We fine-tuned a RoBERTa Base \cite{roberta} classifier on the BAD training set's annotations as the label function $c$ (see Appendix~\ref{subsec:classifier}).
We use the non-toxic part of training data as $\mathcal{D}_\mathrm{LM}$ and all the prompts in the training set as $\mathcal{D}_\mathrm{Pmt}$.
For each prompt $x$, we constructed at most $k=5$ pairs of contrastive samples for $\mathcal{D}_\mathrm{CL}$ from 20 sampled continuations (nucleus sampling, $p=0.9$). We set $\alpha = 0.5$ and $\gamma = 20$ for \textsc{Click}.
For evaluation, all the models generate continuations (response) given the prompts (dialogue history), using nucleus sampling \cite{topp} with $p=0.9$. We conducted simple grid searches for hyperparameters of \textsc{Click} and baselines and selected final values based on the performance on the validation set (see Appendix~\ref{subsec:hyperparameter} for details). See Appendix~\ref{subsec:implementation} for further implementation details.

\paragraph{Baselines}
Following previous work \cite{director, cringe}, we use BlenderBot 365M \cite{blenderbot} as the base model. We compare the following methods.
\textbf{Non-toxic FT} fine-tunes BlenderBot on the non-toxic training set.
\textbf{Unlikelihood} Training \cite{unlikelihood} minimizes the likelihood of each token given the prefix of the toxic sample and also performs language modeling on non-toxic samples.
\textbf{GeDi} \cite{gedi} and \textbf{DExperts} \cite{dexperts} both train a toxic/non-toxic model on the toxic/non-toxic training set and adjust the next-token prediction distribution of the original language model.
\textbf{Director} \cite{director} trains a classification head to similarly adjust the next-token prediction distribution.
\textbf{Cringe} \cite{cringe} improves Unlikelihood Training by applying token-level contrastive learning to toxic samples.

\paragraph{Evaluation Setups} For automatic evaluation, we follow the evaluation metrics in \cite{dexperts}, including the aspects of toxicity, fluency, and diversity. Toxicity is measured by the empirical probability (\textbf{Prob.}) of generating at least one toxic continuation over 25 continuations (labeled by the BAD classifier). Fluency is measured by the mean perplexity (\textbf{Out. PPL}) of generated continuations, as evaluated by a larger language model BlenderBot 1.4B. Diversity is measured using the mean number of distinct $n$-grams, normalized by the text length \cite{distinct}, among the 25 generations for each prompt. We report \textbf{Dist-2/3} scores for distinct bigrams/trigrams, respectively.sa

We also conducted pairwise human evaluation to compare generation results from \textsc{Click} to baselines. 100 prompts were randomly sampled from the BAD test set and each comparison (\textsc{Click} vs. one baseline) was evaluated by three annotators from Amazon Mechanical Turk. Following \cite{dexperts}, evaluation metrics include the perceived level of \textbf{toxicity} (which one is less offensive or biased), \textbf{fluency} (which one is more grammatically correct and coherent), and \textbf{topicality} (which one is more natural, relevant, and logical). See Appendix~\ref{subsec:human-toxicity} for human evaluation details.

\paragraph{Results}
As shown in Table~\ref{tab:toxicity-auto}, \textsc{Click} substantially reduces toxic generations compared to baselines while maintaining reasonable generation diversity. Director and GeDi perform next best to \textsc{Click} but obtain much lower Dist-2/3, indicating that the former two methods both sacrifice generation diversity largely.
Table~\ref{tab:toxicity-human} also shows that human annotators rated \textsc{Click} generations as less toxic than the competitors, demonstrating the effectiveness of \textsc{Click} in eliminating toxic language. See Appendix~\ref{sec:qualitative} for additional qualitative results.

\subsection{Sentiment Steering}
\label{subsec:sentiment}

The task of sentiment steering aims to control the sentiment polarity of generated text, which is well-studied in research of controllable text generation.

\paragraph{Experimental Setups} We evaluated on the test data from \cite{dexperts}, which contains 2.5K/2.5K/5K positive/negative/neutral prompts from OpenWebText \cite{openwebtext}. We use neutral and negative prompts for positive sentiment steering evaluation, and vice versa. The models should generate continuations with either positive or negative sentiment even given prompts with the opposite sentiment (negative or positive, respectively). As in \cite{dexperts}, we use the HuggingFace sentiment classifier as the label function $c$ (see Appendix~\ref{subsec:classifier} for details). For training data, we follow \cite{dexperts} and use SST-5 \cite{sst}. We use sentences with the target sentiment as $\mathcal{D}_\mathrm{LM}$ and the first 2 tokens of all the positive and negative sentences as $\mathcal{D}_\mathrm{Pmt}$.
For each prompt $x$, we constructed at most $k=5$ pairs of contrastive samples for $\mathcal{D}_\mathrm{CL}$ from 20 sampled continuations (nucleus sampling, $p=0.9$) . We set $\alpha=0.1$ and $\gamma=15$ for \textsc{Click}.
For evaluation, all the models generate continuations with maximum 20 tokens using nucleus sampling with $p=0.9$.

\paragraph{Baselines}
We use GPT-2 Large 774M as the base model, consistent with previous work \cite{dexperts}. Same as \S~\ref{subsec:toxicity}, we use \textbf{Target FT}, which fine-tunes GPT-2 on the training data with the target sentiment, \textbf{GeDi}, and \textbf{DExperts} as baselines. We also include \textbf{PPLM} \cite{pplm}, \textbf{CTRL} \cite{ctrl}, and \textbf{DAPT} \cite{dexperts} as baselines. For former two are classical methods for controllable text generation and the latter one applies domain-adaptive pre-training on positive or negative sample corpora. We use these baseline results from \cite{dexperts}.

\paragraph{Evaluation Setups} Following \cite{dexperts}, we report the mean proportion of positive/negative continuations over 25 generated continuations (\textbf{\% Positive/Negative}), as labeled by the HuggingFace sentiment classifier.
\textbf{Out. PPL} is calculated with a larger language model GPT-2 XL 1.5B.
\textbf{Dist-2/3} is calculated consistently with \S~\ref{subsec:toxicity}.

We also conducted pairwise human evaluation for both positive and negative sentiment steering on negative and positive prompts, respectively. Same as \S~\ref{subsec:toxicity}, 100 negative/positive prompts were randomly sampled and each comparison (\textsc{Click} vs. one baseline) was evaluated by three human annotators from the aspects of \textbf{sentiment} (which one is more positive/negative), \textbf{fluency}, and \textbf{topicality}.
See Appendix~\ref{subsec:human-sentiment} for human evaluation details.

\paragraph{Results}
As shown in Table~\ref{tab:sentiment-auto}, \textsc{Click} more effectively steers toward the target sentiments, especially in the adversarial settings (i.e., steering toward the opposite sentiment to the prompt). While \textsc{Click}'s Out. PPL is a bit higher, we believe it is a trade-off with sentiment control since steering a positive/negative prompt toward negativity/positivity may result in an unexpected continuation, which is reflected in a higher Out. PPL.
Table~\ref{tab:sentiment-human} shows that \textsc{Click} has close fluency and topicality to baselines but performs better in sentiment steering. See Appendix~\ref{sec:qualitative} for additional qualitative results.

\subsection{Repetition Reduction}
\label{subsec:repetition}

Autoregressive language models usually suffer from generating text with unnatural repetition \cite{topp}, which is a long-standing and important problem in NLG research \cite{unlikelihood, contrastive-token}. We aim to reduce repetition in language generation with \textsc{Click}.

\paragraph{Experimental Setups} Following previous work \cite{simctg, quark}, we evaluated on the WikiText-103 \cite{wikitext103} dataset, which contains 100M English tokens from Wikipedia articles. We use the official data split as in \cite{unlikelihood, simctg}. We use the diversity metric as the label function $c$, defined as $c(y) = 1 \ \text{if}\ \text{Div}(y) > s \ \text{else}\ 0$, where $\text{Div}(y) = \prod_{n=2}^4 (1.0 - \text{Rep-}n(y)/100)$.
We set $s$ to 0.75, which is the 5\% quantile calculated on human-written text in the training set. We use the WikiText-103 training set as $\mathcal{D}_\mathrm{LM}$ and the first 32 tokens of samples in $\mathcal{D}_\mathrm{LM}$ as $\mathcal{D}_\mathrm{Pmt}$.
For each prompt $x$, we sampled 3/4/5 continuations with $p=0.5/0.7/0.9$, respectively (12 in total, and a lower $p$ usually leads to more repetition), and constructed at most $k=3$ pairs of contrastive samples for $\mathcal{D}_\mathrm{CL}$.
We set $\alpha = 0.3$ and $\gamma = 15$ for \textsc{Click}.
For evaluation, all the models generate continuations with maximum 128 tokens given the prompts with 32 tokens, using greedy decoding where text repetition tends to appear most frequently.

\paragraph{Baselines}
As in \cite{simctg, quark}, we use GPT-2 Base 124M \cite{gpt2} as the base model. We compare \textbf{MLE} (maximum likelihood estimation), the standard language modeling method with the conventional negative-log likelihood loss, \textbf{Unlikelihood} \cite{unlikelihood}, \textbf{SimCTG} \cite{simctg}, a contrastive training method, and \textbf{Quark} \cite{quark}, which conditions language generation on quantized reward tokens. Note that SimCTG, Quark, and our \textsc{Click} are all first pre-trained on the WikiText-103 training set with the MLE objective, and then trained with their own objectives.

\paragraph{Evaluation Setups} We evaluate both the language modeling quality and the generation quality, following previous work \cite{unlikelihood, simctg}. For language modeling quality, we calculate perplexity (\textbf{PPL}) and next-token prediction accuracy (\textbf{Acc}) on the ground-truth continuations of the WikiText-103 test set. We also calculate prediction repetition (\textbf{Rep}), which is defined as the fraction of the next token repeating the prefix tokens, and its variant (\textbf{WRep}), which excludes the cases of the ground-truth token being predicted and repeating the prefix tokens. For generation quality, we report the proportion of repeated 2/3-grams (\textbf{Rep-2/3}) and diversity (\textbf{Div}) as an overall assessment of text repetition. We also report \textbf{MAUVE} \cite{mauve}, an automatic metric that meauses how much the distribution of generated text diverges from human-written text.

We also conducted pairwise human evaluation. 100 prompts were randomly sampled and each pair of generations were compared by three human annotators from the aspects of \textbf{coherence} (which one is more aligned in meaning/topic with the prompt), \textbf{fluency} (which one is more grammatical, understandable, and non-repetitive) and \textbf{overall} quality. See Appendix~\S~\ref{subsec:human-repetition} for human evaluation details.

\paragraph{Results}
As shown in Table~\ref{tab:repetition-auto}, \textsc{Click} remarkably reduces generation repetition with greedy decoding, leading to the highest diversity (0.72) and MAUVE (0.93) scores. While \textsc{Click} has higher PPL and lower Acc, this is probably due to the increased entropy of next-token prediction, which may be a side-product of reducing generation repetition by directly optimizing sequence likelihood. From Table~\ref{tab:repetition-human}, \textsc{Click} is preferred by human in terms of coherence, fluence, and overall quality. See Appendix~\ref{sec:qualitative} for additional qualitative results.

\subsection{Ablation Analysis}
\label{subsec:ablation}

We conduct ablation analysis to give further insights about \textsc{Click}.
We focus on the language detoxification task (\S~\ref{subsec:toxicity}) unless otherwise stated.

\paragraph{Effect of Sample Construction Strategy} We compare \textsc{Click} with several alternatives. For each negative sample $\hat{y}^{-}$,
\textbf{Random} randomly selects a positive sample: $\hat{y}^{+} \in \widehat{\mathcal{Y}}^{+}$, as adopted in previous work \cite{brio, slic},
\textbf{Lower} randomly selects a positive sample only from those with lower likelihood than $\hat{y}^{-}$: $\hat{y}^{+} \in \left\{ \hat{y}^{+} \in \widehat{\mathcal{Y}}^{+} | P_{\theta} (\hat{y}^{+} | x) < P_{\theta} (\hat{y}^{-} | x) \right\}$, 
and \textbf{Lowest} selects the positive sample with the lowest likelihood: $\mathop{\arg \min}_{\widehat{\mathcal{Y}}^{+}} P_{\theta} (\hat{y}^{+} | x)$.

As shown Table~\ref{tab:toxicity-auto}, \ref{tab:sentiment-auto} and \ref{tab:repetition-auto}, \textsc{Click} generally outperforms all the three alternative strategies in either fluency or control effect. We notice that Lower achieves better fluency than Random (lower Out. PPL) in Table~\ref{tab:sentiment-auto}, probably because the former avoids overfitting high-likelihood positive samples. However, Lower and Lowest both underperform \textsc{Click} in fluency (higher Out. PPL in Table~\ref{tab:toxicity-auto} and \ref{tab:sentiment-auto}) and control effect (all the three tables). It confirms our intuitions in \S~\ref{subsec:construction} that exploiting the positive samples with much lower likelihoods than the negative ones somewhat impairs the effectiveness of contrastive learning (biased by contrastive samples with too large likelihood gaps) and the language generation capability (impacted by the low-quality positive samples).

\paragraph{Effect of Weight $\bm{\alpha}$ and Margin $\bm{\gamma}$}
The weight $\alpha$ (Equation \ref{equ:all}) controls the importance of the contrastive loss $\mathcal{L}_\mathrm{CL}$, while the margin $\gamma$ (Equation \ref{equ:cl}) controls the strength of contrastive learning. As shown in Figure~\ref{fig:alpha_gamma}, increasing $\alpha$ and $\gamma$ both lead to lower toxicity (or better controllability), which however sacrifice a bit generation fluency (sightly higher Out. PPL). We speculate this is due to the trade-off between decreasing the generation probability of negative samples (Equation \ref{equ:cl}) and maintaining the underlying language generation capability (Equation \ref{equ:lm}).

\paragraph{Effect of Contrastive Sample Number} In \S~\ref{subsec:toxicity}, we constructed at most $k=5$ pairs of contrastive samples for each prompt $x$. We now vary $k$ from $1$ to $5$. As shown in Figure~\ref{fig:sample_num}, increasing $k$ generally does not reduce toxicity better but instead decreases Out. PPL. The former is probably due to that the contrastive loss (Equation \ref{equ:cl}) has been effective enough to eliminate toxicity. For the latter, we speculate this is because the model-generated positive samples are overall of high likelihood and preferred by the language model (as a reference, the base model BlenderBot generates only 5 toxic ones out of 20 continuations on the BAD training set). Hence, optimization toward more positive samples leads to more generations with similarly high likelihood (or low Out. PPL), as observed in previous work \cite{wang2022exploring}.

\paragraph{Effect of Iterative Training} Similar to the practice in recent work \cite{quark, cringe}, \textsc{Click} can also continue to improve by iterative training (i.e., we use trained \textsc{Click} as the initial model for another iteration). As shown in Table~\ref{tab:toxicity-iter}, \textsc{Click} trained with one additional iteration further reduces toxicity while generation fluency and diversity is slightly impaired. We conjecture it is a trade-off between language generation quality and toxicity, as similarly observed in \cite{quark}.

\section{Related Work}
\label{sec:related}

\paragraph{Controllable Text Generation} As pre-trained language models display the impressive capability of language generation \cite{gpt3}, controlling their generation has become increasingly important in recent years. There are two major directions for controllable text generation: decoding-time and training-based methods.

Decoding-time methods steer model generation toward the desired attribute with lightweight modules without tuning the original model. PPLM \cite{pplm} updates the decoded hidden state according to the classifier's gradient. FUDGE \cite{fudge} trains a classifier to predict whether a partial sequence will satisfy the desired attribute in the future. GeDi \cite{gedi} and DExperts \cite{dexperts} adjust the next-token prediction distribution with two class-conditional auxiliary models. However, decoding-time methods may suffer from high computational expense during generation (e.g., PPLM) and make models inconvenient for out-of-the-box use.

\textsc{Click} falls into training-based methods, which directly train language models to avoid undesirable attributes. Training-based methods include Unlikelihood Training \cite{unlikelihood}, the Cringe loss \cite{cringe}, Quark \cite{quark}, and Director \cite{director}, which are used as compared baselines in our main experiments.

\paragraph{Contrastive Learning for Language Generation} Contrastive learning aims to learn meaningful representations by contrasting positive and negative samples \cite{chen2020simple, he2020momentum, simcse}, which also inspires recent NLG research. CoNT \cite{cont} aligns encoder and decoder representations for non-open-ended language generation. SimCTG \cite{simctg} designs a contrastive training method to learn discriminative and isotropic representations for language generation models. BRIO \cite{brio} and SLiC \cite{slic} uses a contrastive loss to align sequence likelihood with the similarity to reference text. Unlike them, our work applies the contrastive loss to sequence likelihood and targets open-ended text generation tasks, which require special design of sample construction, as discussed in \S~\ref{subsec:prior} and \ref{subsec:ablation}.

\section{Conclusion}

This work introduces a controllable text generation method \textsc{Click}, which needs no modification to the model architecture and facilitates out-of-the-box use of trained models. It employs a contrastive loss on sequence likelihood and adopts a likelihood ranking-based strategy for contrastive sample construction. Our empirical evaluation on the tasks of language detoxification, sentiment steering, and repetition reduction demonstrates that \textsc{Click} can effectively avoid undesirable attributes in language generation and outperforms strong baselines. Ablation analysis gives further insights about \textsc{Click}'s sample construction strategy, hyperparameters, and combination with iterative training. Future work can investigate the combination of \textsc{Click} and various label (or reward) functions \cite{instructgpt}.

\section*{Limitations}
\label{sec:limitation}

Like other controllable text generation methods \cite{pplm, gedi, dexperts, quark, director, cringe}, \textsc{Click} also relies on automatic neural classifiers when constructing $\mathcal{D}_\mathrm{CL}$ in some tasks (language detoxification in \S~\ref{subsec:toxicity} and sentiment steering in \S~\ref{subsec:sentiment} in our work). It may unavoidably inherit the biases and limitations of these classifiers. For instance, for the task of language detoxification, the toxicity may be overestimated when the input prompt or the continuation contains minority identity mentions. To address this limitation, we conducted human evaluation for all the tasks, which further confirms the effectiveness of \textsc{Click}.
As more accurate, inclusive, and reliable classifiers are built (e.g., for toxicity detection), we expect that \textsc{Click} would inherit those improvements as well.

\section*{Ethical Considerations}

As with any controllable text generation technique, \textsc{Click} runs the risk of dual use \cite{pandya2019dual}. Specifically, they could be used to automatically produce harmful contents or malicious behaviors \cite{mcguffie2020radicalization}. Please refer to \cite{bender2021dangers} for a broader discussion of such risks. We hope those who use controllable text generation technologies in real-world deployed systems to consider the potential negative impact and avoid using them to generate harmful contents and misinformation, etc.

For human evaluation, we have obtained study approval from the Institutional Review Board (IRB). We paid the crowdworkers at a fair hourly wage (about \$8/hour) and did not collect any personal identifying information.

\end{document}
